% DOC TYPE E DEF QB SOFTWARE %%%%%%%%%%%%%%%%%%%%%%%%%
\documentclass[12pt]{article}
%%%%%%%%%%%%%%%%%%%% PACKAGES %%%%%%%%%%%%%%%%%%%%
\usepackage[english,italian]{babel}
\usepackage[a4paper, margin=3cm]{geometry}
\usepackage[T1]{fontenc}
\usepackage[utf8]{inputenc}
\usepackage[table]{xcolor}
\usepackage{amsmath}
\usepackage{graphicx}
\usepackage{float}
\usepackage{tabularx}
\usepackage{booktabs}
\usepackage{hyperref}
\usepackage{xcolor}
\usepackage{multicol}
\usepackage{multirow}
\usepackage{soul}
\usepackage{enumitem}
\usepackage{textcomp}
\usepackage{eurosym}
\usepackage{lastpage}
\usepackage{fancyhdr}
\usepackage{adjustbox}
\usepackage{subfiles}

%%%%%%%%%%%%%%%%%%%% COMMAND -> New commands %%%%%%%%%%%%%%%%%%%%
\newcommand{\mailtoQBS}
{
	\href{mailto:qbsoftware.swe@gmail.com}{qbsoftware.swe@gmail.com}
}

\let\oldpar\paragraph
\renewcommand{\paragraph}[1]{\oldpar{#1}\mbox{}\\}

%%%%%%%%%%%%%%%%%%%% COMMAND -> Redefine LaTeX Commands %%%%%%%%%%%%%%%%%%%%
\def\title#1{\gdef\THETITLE{#1}}
\def\date#1{\gdef\THEDATE{#1}}
\def\footname#1{\gdef\THEFOOTNAME{#1}}

\let\oldtexteuro\texteuro
\renewcommand{\texteuro}{\euro}

%%%%%%%%%%%%%%%%%%%% STYLES -> Commands %%%%%%%%%%%%%%%%%%%%
\newcommand{\makefirstpage}
{
	\begin{titlepage}
		
		% Defines a new command for the horizontal lines, change thickness here
		\newcommand{\HRule}{\rule{\linewidth}{0.2mm}} 
		
		\center % Center everything on the page
		
		
		%	Heading Sections
		
		\textsc{\LARGE QB Software}\\[.1cm] 
		\includegraphics[scale=.15]{imgs/qb-software-logo.png}\\[-.1 cm] 
		\includegraphics[scale=.025]{imgs/x.png}\\[0.5cm]
		\includegraphics[scale=.3]{imgs/unipd_logo.png}\\[.5cm]
		\textsc{\Large Università degli studi di Padova}\\[0.5cm] 
		\textsc{\large corso di ingegneria del software }\\[0.5cm]
		\textsc{\large anno accademico 2023/2024 }\\[0.5cm]
		
		
		%	Title section and date section
		
		\ifdefined\THEDATE
		\HRule \\[0.4cm]
		{ \huge{ \bfseries {\THETITLE}} \\ [.5cm]
			\THEDATE}\\[0.4cm] 
		\HRule \\[0.4cm]
		\else
		\HRule \\[0.4cm]
		{ \huge{ \bfseries {\THETITLE}} } \\
		\HRule \\[0.4cm]
		\fi
		%
		\textsc{Contatti:} \mailtoQBS\\[0.3cm]
		
		\vfill 
		
	\end{titlepage}
}

%%%%%%%%%%%%%%%%%%%% ACCESSIBILITY %%%%%%%%%%%%%%%%%%%%
\let\oldhref\href
\renewcommand{\href}[2]{\oldhref{#1}{\textcolor{blue}{\ul{#2}}}}

\hypersetup{
	colorlinks = true,
	linkcolor = cyan,
}

%%%%%%%%%%%%%%%%%%%% STYLES -> Settings %%%%%%%%%%%%%%%%%%%%
% Enable header and footer style
\pagestyle{fancy}
\thispagestyle{empty}
\thispagestyle{fancy}
\pagestyle{fancy}

% Header
\setlength{\topmargin}{-40pt}
\setlength{\headsep}{60pt}
\fancyhf{}
\lhead{QB Software}
\setlength{\headheight}{15pt}
\rhead{\includegraphics[width=1cm]{imgs/qb-software-logo.png}}

% Footer 
\fancyfoot{}
\fancyfoot[L]{\THEFOOTNAME}
\fancyfoot[R]{Pagina \thepage~di~\pageref{LastPage}}
\futurelet\TMPfootrule\def\footrule{\TMPfootrule}
\setcounter{page}{0}
\pagenumbering{arabic}
\renewcommand{\footrulewidth}{0.3pt}

%%%%%%%%%%%%%%%%%%%% ENVIRONMENT %%%%%%%%%%%%%%%%%%%%
% Remove colorless padding in booktabs
\setlength{\aboverulesep}{0cm}
\setlength{\belowrulesep}{0cm}
\setlength{\extrarowheight}{.75ex}

\newenvironment{todo}{
	\rowcolors{2}{cyan!80!black!30!}{cyan!80!black!20!}
	\begin{tabular}{p{3.48cm}>{\raggedright\arraybackslash}p{4cm}>{\raggedright\arraybackslash}p{6.5cm}}
		\toprule
		\rowcolor{gray!20} \textbf{ID}	& \textbf{Interessato} & \textbf{Task} 
		\\\midrule
	}{
		\bottomrule
	\end{tabular}
}

\newenvironment{changelog}{
	\noindent
	{\Large \textbf{Registro delle modifiche}}
	\noindent
	\begin{table}[h]
		\rowcolors{2}{cyan!80!black!30!}{cyan!80!black!20!}
		\begin{adjustbox}{width=\textwidth}
			\begin{tabular}{|c|c|p{2.35cm}|c|p{3.3cm}|}
				\hline
				\rowcolor{gray!20}
				\textbf{V.} & \textbf{Data} & \textbf{Membro} & \textbf{Ruolo} & \textbf{Descrizione} \\
				
				\hline
			}{
				\hline
			\end{tabular}
		\end{adjustbox}
	\end{table}
	
	\clearpage
}

% 1   2    3        4             5            6            7          
% Ver Data Relatore RuoloRelatore DataVerifica Verificatore Descrizione
\newcommand{\newlog}[7]{
	#1 & #5 & #6 & Verificatore & Controllo qualità \\
	   & #2 & #3 & #4           & #7 \\\hline
}

% INFORMAZIONI DOCUMENTO %%%%%%%%%%%%%%%%%%%%%%%%%
\title{Verbale esterno}
\date{del 07 Dicembre 2023}
\footname{Verbale esterno del 07/12/2023}

\begin{document}
	% PRIME PAGINE %%%%%%%%%%%%%%%%%%%%%%%%%
	\makefirstpage
	
	% 1   2    3        4             5            6            7          
% Ver Data Relatore RuoloRelatore DataVerifica Verificatore Descrizione
\begin{changelog}
	\newlog{0.3.0}{15/12/2023}{S. Rovea}{Amministratore}{/12/2023}{}{Correzione ed espansione sezione \ref{sec:management_process}}
	\newlog{0.2.0}{02/12/2023}{S. Destro}{Amministratore}{03/12/2023}{A. Feltrin}{Aggiunta sezione su processi organizzativi, sezione \ref{sec:management_process}}
	\newlog{0.1.0}{22/11/2023}{A. Bustreo}{Amministratore}{22/11/2023}{A. Giurisato}{Aggiunto processo supporto per la documentazione, sezione \ref{sec:documentation}. Aggiunta prefazione, sezione \ref{sec:prefazione}}
\end{changelog}
	\clearpage
	
	\tableofcontents
	\clearpage

    \section{Informazioni generali}
    
    \subsection{Luogo e data dell'incontro}
    
    \begin{itemize}
    	\item \textbf{Luogo}: Comunicazione su Discord
    	\item \textbf{Data}: 07/12/2023
    	\item \textbf{Ora di inizio}: 14:00
    	\item \textbf{Ora di fine}: 18:00
    \end{itemize}
    
    \subsection{Presenze}
    
    \begin{itemize}
    	\item \textbf{Totale presenze}:
    	\begin{itemize}
    		\item Bustreo Alessandro
    		\item Destro Stefano
    		\item Domuta Alessia 
    		\item Feltrin Alessandro 
    		\item Fontana Raffaele Paolo 
    		\item Giurisato Andrea 
    		\item Rovea Silvia
    	\end{itemize}
    	
    	\item \textbf{Totale assenze}: 0
    	
    	\item \textbf{Partecipanti esterni}:
    	\begin{itemize}
    		\item Federico Rispo (Zextras)
    	\end{itemize}
    \end{itemize}
    
    \section{Ordine del giorno}
    \begin{itemize}
    	\item domande da porre all'azienda.
    \end{itemize}
    
    \section{Verbale}
    Le domande poste da QB Software a Zextras durante la conversazione:
    
    % Risposta
    \newcommand{\answer}{\item[\textbf{A:}]}

    \begin{enumerate}[label=\textbf{Q\arabic*:}]
        \item Dobbiamo considerare anche l'inserimento di allegati da parte di un utente oppure non è necessario per il progetto?
        \answer Per quanto riguarda gli allegati non è richiesta alcun tipo di funzionalità quindi ignorate pure tale casistica.

        \item Avevamo dei dubbi riguardanti il concetto di oggetto presente nel capitolato a pagina 4 nella sezione Requisiti Funzionali, c'è scritto "l’eliminazione di un oggetto": ci si riferisce soltanto alle email o anche ad altro?
        \answer Confermo che con oggetto si intende un'email (un contatto o un evento a calendario se decidete di implementare anche i requisiti funzionali opzionali riguardo questi altri due domini). Se implementerete solo il dominio delle email  allora la cancellazione sarà solo delle email.

        \item Per quanto riguarda le cartelle non ci è chiaro se stiamo parlando di cartelle come posta arrivata/posta inviata, o se stiamo parlando di una cartella che contiene file?
        \answer Con cartelle si intende la creazione di una cartella che possa contenere email. L'esempio delle cartelle come posta arrivata/posta inviata è corretto. La creazione di cartelle che contengono email è un concetto che è presente anche in sistemi come Google Workspace o Office 365.

        \item Non ci è familiare il concetto di condividere una cartella come posta in arrivo o posta inviata. Abbiamo capito bene il requisito di condivisione?
        \answer Avete capito bene. La condivisione di tali cartelle è una funzionalità richiesta soprattutto in ambienti aziendali quando si vuole condividere con i colleghi email che riceve una singola persona.  Avendo una cartella condivisa ci si può spostare dentro tutte quelle email che si vogliono condividere con le persone a cui è stato dato accesso di lettura alla cartella.

        \item Vorremmo sapere nello specifico gli stress test da eseguire nel caso in cui decidessimo di implementare quel requisito opzionale. Quale tipologia di test dobbiamo eseguire? 
        \answer Sarebbe meglio avere uno stress test per ogni requisito: per esempio l'invio di email.
    \end{enumerate}
    
    \section{Azioni da intraprendere}
        \begin{itemize}
            \item Continuare con la stesura dei Casi d'Uso.
    	\end{itemize}
\end{document}